% ------------------------------------------------------------
% Chapter 5: Redshift Without Metric Expansion
% ------------------------------------------------------------

\chapter{Redshift Without Metric Expansion}
\label{chap:redshift}

\section{Overview}

A distinctive feature of the \rsvpabbr{} cosmology is the possibility of
explaining observed redshift entirely through lamphrodynamic effects rather
than metric expansion. Whereas standard cosmology attributes the redshift of
light to the stretching of spacetime (increase of the scale factor $a(t)$),
\rsvpabbr{} allows for redshift generated by entropy gradients and the
interaction of null geodesics with the lamphrodynamic fields
$(\PhiField,\vField,\SField)$.

In this chapter we provide a complete derivation of the redshift formula
\[
  z = \int_\gamma f(\SField,\vField)\,d\tau,
\]
starting from the modified Einstein equations and the null geodesic
equations from \cref{chap:geodesics}. We then compare the resulting
distance--redshift relation with the standard $z \approx H_0 d/c$ law and
derive higher--order corrections.

\section{Null Geodesics in Lamphrodynamic Backgrounds}

Consider a null curve $\gamma(\tau)$ parameterized by affine parameter $\tau$
satisfying the modified geodesic equation
\begin{equation}
\label{eq:null-geodesic}
  \frac{d^2 x^\lambda}{d\tau^2}
  + \Gamma^\lambda_{\mu\nu}
  \frac{dx^\mu}{d\tau}\frac{dx^\nu}{d\tau}
  + \Delta^\lambda_{\mu\nu}
  \frac{dx^\mu}{d\tau}\frac{dx^\nu}{d\tau}
  = 0,
\end{equation}
with $\Gamma$ the Levi--Civita connection and $\Delta$ given by
\eqref{eq:delta-connection}. For null curves,
\[
  g_{\mu\nu}\frac{dx^\mu}{d\tau}\frac{dx^\nu}{d\tau}=0.
\]

Let $k^\mu = dx^\mu/d\tau$ denote the null tangent vector. Contracting
\eqref{eq:null-geodesic} with $g_{\lambda\rho}$ yields
\[
  k^\nu \nabla_\nu k_\rho
  = -\Delta_{\rho\mu\nu} k^\mu k^\nu,
\]
where $\Delta_{\rho\mu\nu} := g_{\rho\lambda}\Delta^\lambda_{\mu\nu}$.
Thus, lamphrodynamics introduces an effective \emph{drag} or
\emph{frequency--shift} term along null rays.

\section{Frequency Transport Equation}

Let $\omega(\tau)$ denote the photon frequency measured by a comoving
observer with 4--velocity $u^\mu$. Then
\[
  \omega = - g_{\mu\nu} k^\mu u^\nu.
\]
Differentiating along $\gamma$, we obtain
\begin{equation}
\label{eq:omega-transport}
  \frac{d\omega}{d\tau}
  = -u^\nu k^\mu k^\rho \nabla_\rho g_{\mu\nu}
  - g_{\mu\nu} k^\mu k^\rho \nabla_\rho u^\nu
  - g_{\mu\nu} k^\mu k^\rho \Delta^\nu_{\rho\sigma} u^\sigma.
\end{equation}
For a comoving observer in an FLRW background with lamphrodynamics,
the first two terms reproduce the standard cosmic redshift associated
with the scale factor, while the last term introduces new entropy--flux
effects.

Using the homogeneous ansatz from \cref{sec:flrw-ansatz} and the expression
for $\Delta^\nu_{\rho\sigma}$ from \eqref{eq:delta-connection},
\eqref{eq:omega-transport} simplifies to
\begin{equation}
\label{eq:omega-transport2}
  \frac{1}{\omega}\frac{d\omega}{d\tau}
  = -\frac{\dot{a}}{a}
  - \alpha\,k^\mu k^\rho \nabla_\rho(\SField\,\vField_\mu)
  + \mathcal{O}(\alpha^2).
\end{equation}

\section{Case I: No Metric Expansion}

One may consistently set $a(t)\equiv 1$ (or constant) if spatial sections of
the plenum are stationary in the lamphrodynamic gauge. Then the standard
cosmic redshift term disappears, leaving
\begin{equation}
\label{eq:no-expansion}
  \frac{1}{\omega}\frac{d\omega}{d\tau}
  = - \alpha\,k^\mu k^\rho \nabla_\rho(\SField\,\vField_\mu).
\end{equation}
Integrating along the null geodesic from emission $(e)$ to observation $(o)$
gives
\begin{equation}
\label{eq:z-integral}
  1+z
  = \frac{\omega(e)}{\omega(o)}
  = \exp\left(
      \alpha\int_{\gamma_{e\to o}}
      k^\mu k^\rho \nabla_\rho(\SField\,\vField_\mu)\,d\tau
    \right).
\end{equation}
Thus,
\[
  z = \exp\left(
      \alpha\int_{\gamma} k^\mu k^\rho \nabla_\rho(\SField\,\vField_\mu)\,d\tau
    \right) - 1.
\]

Defining the lamphrodynamic integrand
\begin{equation}
\label{eq:f-definition}
  f(\SField,\vField)
  := \alpha\,k^\mu k^\rho \nabla_\rho(\SField\,\vField_\mu),
\end{equation}
we recover the expression used in heuristic discussions of RSVP cosmology:
\[
  z = \int_\gamma f(\SField,\vField)\,d\tau
\]
to lowest order in $f$.

\section{Comparison with Standard Hubble Law}

For small $z$,
\[
  z
  \approx \alpha \int_\gamma k^\mu k^\rho
               \nabla_\rho(\SField\,\vField_\mu)\,d\tau.
\]
Let $d$ denote the comoving distance along the null ray and assume that
$\SField$ and $\vField$ vary slowly on cosmological timescales. Then
\[
  z \approx \alpha\,(\nabla\SField\cdot\vField)\,d/c,
\]
suggesting an effective Hubble constant
\[
  H_{\mathrm{eff}}
  := \alpha\,(\nabla\SField\cdot\vField).
\]
In this interpretation, cosmic redshift is governed not by $a(t)$ but by
lamphrodynamic entropy flux along null trajectories.

\section{Higher–Order Corrections}

Including higher–order terms in $\Delta^\lambda_{\mu\nu}$ and accounting for
the evolution of $\SField$ and $\vField$ with cosmic time, we obtain
\[
  z = H_{\mathrm{eff}}\frac{d}{c}
  + \gamma_2 \left(\frac{d}{c}\right)^2
  + \gamma_3 \left(\frac{d}{c}\right)^3 + \cdots,
\]
with coefficients $\gamma_i$ determined by lamphrodynamic couplings and
cosmic entropy production history. Unlike $\Lambda$CDM, the expansion
history is replaced by entropy history, leading to distinct observational
signatures (see \cref{chap:supernovae}).

\section{Distance–Redshift Relation}

Inverting $z(d)$ yields a nontrivial distance–redshift relation. For
small $z$,
\[
  d \approx \frac{c}{H_{\mathrm{eff}}} z
    - \frac{\gamma_2 c}{H_{\mathrm{eff}}^3} z^2
    + \mathcal{O}(z^3),
\]
which should be compared with the standard $\Lambda$CDM luminosity
distance–redshift relation. Observational consequences are treated in
detail in \cref{chap:supernovae}.

\section{Discussion}

Lamphrodynamic redshift provides an alternative explanation for cosmological
redshifts that does not require metric expansion. The key observable
predictions lie in deviations from the standard Hubble law at moderate
redshifts and in the detailed luminosity distance relations. In the next
chapter (\cref{chap:structure-formation}), we turn to structure formation
and the implications of lamphrodynamics for the growth of density
perturbations.

\clearpage

